% =====================================================================
\section{领头区域：软区（非超软）、快度与 Glauber 区}\label{sec:tmd-regions}
% =====================================================================

与共线因子化相比，小 $q_T$ 的区域分析有两处本质的不同。其一，测量量 $\mathbf q_T$ 约束所有部分子的横动量之和，因此共线子图的横动量 $k_T\sim q_T$ 不能被积掉，软辐射的横动量也不能被积掉：一个软胶子只要带有 $q_T$ 量级的横动量，就以领头幂次改变测量值。其二，在固定横动量 $k_T$ 下，在壳的无质量动量由快度 $y=\frac12\ln(k^+/k^-)$ 完全确定：$k=(\frac{k_T}{\sqrt2}\ee^{y},\frac{k_T}{\sqrt2}\ee^{-y},\mathbf k_T)$。$A$ 共线的动量 $Q(1,\lambda^2,\lambda)$ 对应 $y\approx\ln(Q/q_T)$，$B$ 共线的 $Q(\lambda^2,1,\lambda)$ 对应 $y\approx-\ln(Q/q_T)$，软的 $Q(\lambda,\lambda,\lambda)$ 对应 $|y|\lesssim1$；三者的虚度都是 $q_T^2$，它们之间是一条连续的快度谱，区域之间只能用快度区分。这正是快度发散的来源：用类光 Wilson 线表示软子图时，沿快度的积分在 $y\to\pm\infty$ 处不受约束。

\begin{proposition}[小 $q_T$ 的领头区域 \tagA]\label{prop:tmd-regions}
对\cref{def:tmd-kin}中的过程，在 Feynman 规范与 $\lambda=q_T/Q$ 的幂次计数下，领头幂次的约化图由以下子图组成：硬子图 $H$（及割线另一侧的 $H^\dagger$），全部内线的虚度 $\sim Q^2$；$A$ 共线子图 $J_A$ 与 $B$ 共线子图 $J_B$，动量标度 $Q(1,\lambda^2,\lambda)$ 与 $Q(\lambda^2,1,\lambda)$，其横动量 $k_T\sim q_T$ 保留在测量函数中；一个\emph{软}子图 $S$，其动量的全部分量都是 $q_T$ 量级，即 $Q(\lambda,\lambda,\lambda)$——而不是\cref{def:ir-scalings}的超软标度 $Q(\lambda^2,\lambda^2,\lambda^2)$——它只经胶子与 $J_A$、$J_B$ 相连，不接入 $H$；以及（对 Drell--Yan）两个强子的旁观者之间的 Glauber 交换，标度 $Q(\lambda^2,\lambda^2,\lambda)$。每个喷注在割线每侧以恰好一个物理部分子加上任意多个纵向（标量极化）共线胶子与 $H$ 相连。软区与共线区只在快度上区分。（本命题只涉及部分子态。）
% errtype: E-S
\uses{thm:ir-leading-regions,prop:ir-power-counting,def:tmd-kin}
\end{proposition}

\begin{proofsketch}
(1) \emph{测量函数。}对 Drell--Yan，$\delta^{(2)}(\mathbf q_T-\sum_{i\in J_A}\mathbf k_{iT}-\sum_{i\in J_B}\mathbf k_{iT}-\sum_{i\in S}\mathbf k_{iT})$ 只约束进入末态的横动量之和；因而对 $q_T\sim\lambda Q$，横动量 $\gg q_T$ 的末态辐射只能来自 $H$（它们给出 $q_T\sim Q$ 的贡献，归入 $Y$ 项与幂次修正），横动量 $\ll q_T$ 的辐射不改变测量值。
(2) \emph{幂次计数。}把\cref{prop:ir-power-counting}的推动—量纲论证用于软标度 $Q(\lambda,\lambda,\lambda)$：一个软胶子连接 $J_A$ 与 $J_B$ 时，测度 $\dd^4k\sim\lambda^4$，胶子传播子 $\sim\lambda^{-2}$，两条被软动量带离壳的共线线各给出一个程函分母 $p_A^+k^-\sim\lambda$，与其分子 $p_A^\mu$ 合起来为 $\lambda^{-1}$，合计 $\lambda^{4-2-1-1}=\lambda^0$，与超软标度给出的 $\lambda^{8-4-2-2}$ 相同；但超软胶子的横动量 $\lambda^2Q\ll q_T$ 不改变测量值，它们是软区的低端，并入软子图。软线接入 $H$、软夸克附着于喷注、喷注以多于一个物理部分子接入 $H$ 都按\cref{prop:ir-power-counting}被幂次压低（此步只给出 E-S 型陈述）。
(3) \emph{区域列表。}由 (1)(2) 与\cref{thm:ir-leading-regions}（全阶论证，\tagA）的钳夹面分类，即得所列的 $H$--$J_A$--$J_B$--$S$ 结构；与共线因子化的差别只在于软子图的标度与测量函数。
(4) \emph{Glauber 区。}旁观者之间的 Glauber 交换按\cref{thm:ir-leading-regions}的第 (4) 步同样是 $\lambda^0$，并在单个割图中被钳住，须单独处理（\cref{lem:tmd-glauber}）。
文献：\cite{CollinsSoper1981}；\cite[第~10~章与第~13.3~节]{Collins2011}。
\end{proofsketch}

\paragraph{地位与缺口}\tagA。论证是标度陈述（E-S）：它确定了每个区域的幂次，但没有给出积分形式的估计；积分形式就是\cref{asm:appN-envelope}与\cref{asm:appN-separation}（均为 \tagA），在 TMD 的快度窗口中还需要\cref{asm:appN-tmd}(T1)，即横动量很小的线的尾部贡献被压低。区域的完备性（没有遗漏的领头区域）继承自\cref{thm:ir-leading-regions}，其中最弱的一步是嵌套区域的幂次计数与 Feynman 规范中额外不连通硬散射的相消\cite[第~14.4.8~节]{Collins2011}。对强子态，区域的完备性正是\cref{asm:intro-bridge}(Br1)。升级途径：在快度窗口内对每个区域证明积分范数下的幂次界（\cref{rem:tmd-rapidity-window}）。

软区与超软区的区别在这里是实质性的。对共线单举的观测量，软胶子只需要满足 $k^\pm\ll Q$，Libby--Sterman 的“软”区（本书的超软区）就够用了；对 $q_T$ 微分的截面，测量量本身的标度是 $q_T=\lambda Q$，与之相应的软标度是 $\lambda Q$ 而不是 $\lambda^2Q$。在 $(\ln k^+,\ln k^-)$ 平面上，$k_T\sim q_T$ 的在壳动量都位于双曲线 $2k^+k^-=q_T^2$ 附近，$A$ 共线、软、$B$ 共线依次沿这条双曲线排列；以类光 Wilson 线代替共线部分子时，软子图对整条双曲线积分，在两端（快度 $\to\pm\infty$）发散——这就是快度发散，它反映的是软区与共线区在快度上的重叠，而不是某个正规子的人为产物。下一节的定义就是为了把这一重叠只计一次。
